\documentclass[11pt]{article}

\usepackage[margin=1in]{geometry}
\usepackage{amsmath,amssymb,amsthm,mathtools}
\usepackage{enumitem}
\usepackage{hyperref}

\setlength{\parskip}{0.5em}
\setlength{\parindent}{0pt}

\newcommand{\R}{\mathbb{R}}
\newcommand{\K}{\mathcal{K}}
\newcommand{\eps}{\varepsilon}
\newcommand{\bstar}{f^\ast}
\newcommand{\xstar}{x^\ast}
\newcommand{\norm}[1]{\lVert #1\rVert}
\newcommand{\E}{\mathbb{E}}
\newcommand{\Tr}{\operatorname{Tr}}
\newcommand{\op}{\mathrm{op}}

\begin{document}

\begin{center}
{\Large \textbf{DSC 243 --- Homework 2}} \\
\vspace{0.5em}
{\large The Positive Semidefinite Case and Convergence Under Spectral Structure}
\end{center}

\textbf{Instructions.}
Use the notation from the notes. Unless explicitly stated otherwise, assume
\begin{align*}
&f(x) = \tfrac{1}{2} x^\top A x - b^\top x,\\
&A = A^\top \succeq 0,\\
&b \in \mathrm{range}(A),\\
&\beta = \lambda_{\max}(A),\\
&\alpha = \lambda_{\min}(A),\\
&\text{and when } \alpha > 0,\ \kappa = \beta/\alpha.
\end{align*}
Let $\xstar\in\{x:Ax=b\}$ and $e_0 = x_0 - \xstar$. Let $v_1,\dots,v_d$ be an orthonormal eigenbasis of $A$ with eigenvalues $\lambda_1,\dots,\lambda_d$, and write $e_0 = \sum_{i=1}^d c_i v_i$. Throughout, we repeatedly use the exact spectral formula
\[
f(x_k) - \bstar \;=\; \tfrac{1}{2}\sum_{i=1}^{d}\lambda_i\,p_k(\lambda_i)^2\,c_i^2, \qquad p_k(\lambda) := \prod_{j=1}^{k}(1-\eta_j\lambda),
\]
derived in Section 2 for gradient descent with time-varying stepsizes $\eta_1,\dots,\eta_k$. For stochastic least squares, use the notation
\[
L(w)=\frac12\E[(y-\langle w,x\rangle)^2],\qquad H=\E[xx^\top],\qquad \|u\|_H^2=u^\top H u.
\]
You may quote results proved in class or in the notes, but every nontrivial step must be justified.

\section*{Part I: Theory}

\begin{enumerate}[label=\textbf{T\arabic*.}, leftmargin=0pt, itemindent=2.5em]

%---------------------------------------------------------------------
\item \textbf{(Optional) Mini-batch SGD and the bias--variance tradeoff.}
Assume the stochastic least-squares setup from the notes: $H=\E[xx^\top]\succ 0$, $H\succeq \mu I$, and
\[
\E[\|x\|^2xx^\top]\preceq R^2H.
\]
Let
\[
\Sigma:=\E[(y-\langle w_\ast,x\rangle)^2xx^\top].
\]
For an integer batch size $B\ge 1$, define mini-batch SGD by
\[
w_t=w_{t-1}+\frac{\gamma}{B}\sum_{j=1}^B
\bigl(y_{t,j}-\langle w_{t-1},x_{t,j}\rangle\bigr)x_{t,j},
\]
where the samples $(x_{t,j},y_{t,j})$ are all independent copies of $(x,y)$.
\begin{enumerate}[label=(\alph*)]
\item Define
\[
\widehat H_t:=\frac1B\sum_{j=1}^B x_{t,j}x_{t,j}^\top,
\qquad
\overline\xi_t:=-\frac1B\sum_{j=1}^B
(y_{t,j}-\langle w_\ast,x_{t,j}\rangle)x_{t,j}.
\]
Show that the error $e_t=w_t-w_\ast$ satisfies
\[
e_t=(I-\gamma\widehat H_t)e_{t-1}-\gamma\overline\xi_t.
\]
\item Prove that
\[
\E[\widehat H_t^2]\preceq R^2H.
\]
\item Show that the bias process satisfies the same contraction as in the single-sample proof:
\[
\E\|b_t\|_H^2\le R^2e^{-\gamma\mu t}\|w_0-w_\ast\|^2,
\qquad 0<\gamma<1/R^2.
\]
\item Show that
\[
\E[\overline\xi_t\overline\xi_t^\top]=\frac{1}{B}\Sigma.
\]
\item Adapt the covariance argument from the notes to prove the last-iterate bound
\[
\E[L(w_t)]-L(w_\ast)
\le
R^2e^{-\gamma\mu t}\|w_0-w_\ast\|^2
+\frac{\gamma\Tr(\Sigma)}{B(2-\gamma R^2)}.
\]
\item Suppose the total sample budget is $N=Bt$. According to the bound in part (e), increasing $B$ lowers the variance floor but reduces the number of parameter updates. Explain the resulting tradeoff qualitatively. In which regime would you prefer $B=1$, and in which regime would a larger batch be attractive?
\end{enumerate}

\noindent
[\textbf{Hint for part (b):} Expand $\widehat H_t^2$. The diagonal terms use the fourth-moment bound, while the off-diagonal terms use independence and the inequality $H^2\preceq \|H\|_{\op}H\preceq R^2H$.]

%---------------------------------------------------------------------
\item \textbf{(Optional) A near-critical Marchenko--Pastur crossover.}
For the Marchenko--Pastur law with aspect ratio $\gamma<1$, the left edge is $\lambda_-=(1-\sqrt\gamma)^2>0$ and the GD spectral integral has an exponential backbone. At $\gamma=1$, the left edge hits zero and the rate becomes purely polynomial.

Suppose that the aspect ratio depends on the iteration horizon through
\[
\sqrt{\gamma_k}=1-\frac{\tau}{\sqrt{k}},
\qquad \tau>0,
\]
so that $\lambda_-\asymp \tau^2/k$. Assume isotropic error weights and use the Marchenko--Pastur density.
\begin{enumerate}[label=(\alph*)]
\item Write the GD spectral integral with $\eta=1/\lambda_+$ and make the change of variables $\lambda=\lambda_-+u/k$.
\item Show formally that the exponential factor no longer behaves like a fixed geometric rate; instead it converges to a nontrivial function depending on $\tau$ and $u$.
\item Use this calculation to explain why the regime $\lambda_-\asymp 1/k$ is a crossover between positive-definite exponential convergence and hard-edge polynomial convergence.
\end{enumerate}

\end{enumerate}

\section*{Part II: Experiments}

\begin{enumerate}[label=\textbf{E\arabic*.}, leftmargin=0pt, itemindent=2.5em]

%---------------------------------------------------------------------
\item \textbf{GD, PSD Chebyshev, and CG on a power-law spectrum.}
Take $d=200$ and
\[
A \;=\; \operatorname{diag}\bigl(i^{-3}\bigr)_{i=1}^d
\]
(so $A\succ 0$ with $\alpha=d^{-3}$, $\beta=1$ and $\kappa=d^3\approx 8\times 10^6$). Let $\xstar=\mathbf{1}\in\R^d$, $b=A\xstar$. Use the same generic random initial point $x_0$ for every method.

\begin{enumerate}[label=(\alph*)]
\item Run gradient descent with constant stepsize $\eta=1/\beta$ for $k=5000$ iterations.

\item For each horizon $K\in\{10,20,50,100,200,500,1000,2000\}$, run PSD-Chebyshev with the stepsize schedule $\eta_j = 1/(\beta\sin^2(j\pi/(2K)))$ for $j=1,\dots,K$, starting from the same $x_0$. Record the final $f(x_K)-\bstar$ for each choice of $K$ and plot these values as a function of $K$.

\item Run the conjugate gradient method from $x_0$ for $k=5000$ iterations (truncating once the gap drops below $10^{-13}\norm{x_0-\xstar}^2$).

\item On the same semilog plot of $(f(x_k)-\bstar)/(f(x_0)-\bstar)$ versus $k$, overlay:
\begin{itemize}
\item the GD curve from (a);
\item the final-iterate values from (b), plotted as a step function or marker plot;
\item the CG curve from (c);
\item the theoretical upper bounds $\beta/(2(2k+1))\cdot \norm{x_0-\xstar}^2/(f(x_0)-\bstar)$ (Theorem 6.1) and $\beta/(8k^2)\cdot \norm{x_0-\xstar}^2/(f(x_0)-\bstar)$ (Theorem 6.2).
\end{itemize}

\item Comment on the empirically observed rates of GD, PSD-Chebyshev, and CG. Which is fastest? For large $k$, does CG separate from PSD-Chebyshev, and if so, why? (Recall that the spectrum here decays polynomially, so the spectral-density analysis of Section~7 is relevant.)
\end{enumerate}

%---------------------------------------------------------------------
\item \textbf{Source condition and the matching GD rate.}
Take $d=2000$ and
\[
A \;=\; \operatorname{diag}\bigl(i^{-\mathsf{a}}\bigr)_{i=1}^d\qquad\text{with }\mathsf{a}=2.
\]

For each $s\in\{-\tfrac{1}{4},\,0,\,\tfrac{1}{2},\,1\}$, generate an initial error satisfying the source condition $e_0=A^s w$ as follows: sample $w\in\R^d$ with iid $\mathcal{N}(0,1)$ entries (rescale so $\norm{w}^2=d$), and set $x_0=\xstar+e_0$ with $\xstar=\mathbf{1}$, $b=A\xstar$.

\begin{enumerate}[label=(\alph*)]
\item For each $s$, run GD with $\eta=1/\beta$ for $k=10^4$ iterations. Plot $(f(x_k)-\bstar)/(f(x_0)-\bstar)$ versus $k$ on a log-log scale (four curves, one per value of $s$).

\item On each curve, overlay the theoretical reference line $c_s\,k^{-(1+2s)}$ (dotted), where $c_s$ is the theoretical constant from Theorem~7.1 (expressed in terms of $\beta$, $s$, $\norm{w}^2$, $f(x_0)-\bstar$). Do the empirical slopes match the theoretical exponent $-(1+2s)$?

\item For $s=-1/4$ (the case with $s<0$), compare the empirical value of $\norm{e_0}^2$ and $f(x_0)-\bstar$. Observe that $\norm{e_0}^2$ is large but the function gap remains well-defined. Explain in one or two sentences how Theorem~7.1 delivers a \emph{dimension-free} bound in this regime (that is, the right-hand side of the bound depends on $\norm{w}$, not $\norm{e_0}$).

\item Compute iteration complexity empirically: for each $s$, report the smallest $k$ that achieves relative gap $10^{-4}$ and $10^{-6}$. Does the ratio of complexities across $s$ match what T3(e) predicts?
\end{enumerate}

%---------------------------------------------------------------------
\item \textbf{CG on Marchenko--Pastur spectra across the three regimes.}
Fix $d=1200$ and consider three aspect ratios $\gamma\in\{0.5,\,1,\,2\}$:
\[
n_{0.5} \;=\; 2d,\qquad n_{1} \;=\; d,\qquad n_{2} \;=\; d/2.
\]
For each $\gamma$, draw $D\in\R^{n\times d}$ with iid $\mathcal{N}(0,1)$ entries and set $A=\tfrac{1}{n}D^\top D$. Sample $\xstar$ uniformly on the sphere of radius $\sqrt{d}$ and set $b=A\xstar$, $x_0=0$. Run both GD with stepsize $\eta=1/\lambda_{\max}(A)$ and CG, and record $f(x_k)-\bstar$ versus $k$. For each $\gamma$, take the \emph{median} over $20$ independent draws of $D$ (and $\xstar$).

\begin{enumerate}[label=(\alph*)]
\item For each $\gamma$, plot on a log-log scale the median relative gap $(f(x_k)-\bstar)/(f(x_0)-\bstar)$ versus $k$ for GD and CG. Shade the $10\%$--$90\%$ interquantile range.

\item On the $\gamma=1$ plot, overlay the dotted reference lines
\begin{align*}
\text{GD reference:}\quad &\frac{1}{\sqrt{2\pi}\,k^{3/2}}\cdot\frac{\norm{e_0}^2}{f(x_0)-\bstar},\\[2pt]
\text{CG reference (CG on MP at $\gamma=1$, see notes):}\quad &\frac{3}{(k+1)(k+2)(2k+3)}\cdot\frac{\norm{e_0}^2}{f(x_0)-\bstar}.
\end{align*}
Report whether the empirical rates track the predicted $k^{-3/2}$ and $k^{-3}$ exponents.

\item On the $\gamma=0.5$ plot (positive definite case), overlay the \emph{theoretical Chebyshev upper bound}
\[
4\left(\frac{\sqrt{\kappa}-1}{\sqrt{\kappa}+1}\right)^{2k}\cdot 1,\qquad \kappa\;=\;\bigl((1+\sqrt\gamma)/(1-\sqrt\gamma)\bigr)^2,
\]
and also a power-times-exponential reference $c\,k^{-3/2}(1-1/\kappa)^{2k}$ (from the Laplace-method Theorem~7.3 applied to the MP edge). Which of the two is closer to empirical CG for early $k$? Late $k$? Briefly explain why.

\item On the $\gamma=2$ plot (rank-deficient), verify that in the effective spectrum the \emph{nonzero} eigenvalues lie in $[(\sqrt\gamma-1)^2,(\sqrt\gamma+1)^2]$. Overlay the PSD Chebyshev/CG worst-case bound $\beta/(8k^2)\cdot\norm{e_0}^2$ as a reference, as well as the power-times-exponential reference from the Laplace method with $\alpha_{\mathrm{eff}}=(\sqrt\gamma-1)^2$ and $\beta=(\sqrt\gamma+1)^2$. Discuss whether the effective-condition-number rate or the power-times-exponential rate better describes empirical behavior.

\item \emph{Summary.} In one or two sentences, describe how the three MP regimes ($\gamma<1$, $\gamma=1$, $\gamma>1$) differ from the point of view of convergence of CG: where does the $k^{-3/2}$ edge factor come from, and in which regime does CG achieve the sharpest gain over GD?
\end{enumerate}

\end{enumerate}

\end{document}
